\documentclass[11pt]{beamer}

\usetheme{Madrid}
\usecolortheme{default}

\usepackage[utf8]{inputenc}
\usepackage[T1]{fontenc}
\usepackage{amsmath, amssymb, physics, bm}

\title[Fractal Electrodynamics]{Electrodynamics of Fractal Distributions of Charges and Fields}
\subtitle{Fractional Continuous Models and Generalized Maxwell Equations}
\author{Based on V.E. Tarasov}
\institute{By Gytis Vėjelis}
\date{}

\begin{document}

%------------------------------------------------
\begin{frame}
  \titlepage
\end{frame}

%------------------------------------------------
\begin{frame}{Motivation}
\begin{itemize}
  \item Classical electrodynamics assumes \textbf{smooth charge and current distributions} in $\mathbb{R}^3$
  \item Many physical systems are irregular or strongly inhomogeneous:
  \begin{itemize}
    \item Turbulent plasmas
    \item Porous or disordered media
    \item Percolation clusters
    \item Space plasmas and astrophysical environments
  \end{itemize}
  \item These systems often exhibit \textbf{fractal geometry} with non-integer dimension
\end{itemize}
\vspace{0.3cm}
How should Maxwell's equations be modified when charges live on fractals?
\end{frame}

%------------------------------------------------
\begin{frame}{What is a Fractal Distribution?}
\begin{itemize}
  \item A fractal set has a (typically non-integer) \textbf{Hausdorff dimension} $D$
  \item Charge does not fill 3D volume uniformly, but resides on a set of dimension $D$
  \item Scaling law for total charge inside a ball of radius $R$:
\end{itemize}
\begin{equation*}
  Q(R) \propto R^D, \qquad 0 < D \le 3
\end{equation*}
\begin{itemize}
  \item In ordinary media, $D = 3$ for a bulk, $D = 2$ for a surface, $D = 1$ for a line
  \item Here $D$ can be non-integer and is defined operationally from the scaling
\end{itemize}
\end{frame}

%------------------------------------------------
\begin{frame}{Standard Electrodynamics}
\begin{itemize}
  \item Charge density $\rho(\bm r)$ defined on $\mathbb{R}^3$
  \item Volume, surface, and line elements:
\end{itemize}
\begin{equation*}
  dV = d^3x, \qquad d\bm S, \qquad d\bm l
\end{equation*}
\begin{itemize}
  \item Total charge in region $W \subset \mathbb{R}^3$:
\end{itemize}
\begin{equation*}
  Q(W) = \int_W \rho(\bm r)\, dV
\end{equation*}
\begin{itemize}
  \item Current density $\bm J(\bm r, t)$ and total current through surface $S$:
\end{itemize}
\begin{equation*}
  I(S) = \int_S \bm J \cdot d\bm S
\end{equation*}
\end{frame}

%------------------------------------------------
\begin{frame}{Standard Electrodynamics}
Maxwell equations (differential form, in vacuum):
\begin{align*}
  \nabla \cdot \bm E &= \frac{\rho}{\varepsilon_0} , \\
  \nabla \cdot \bm B &= 0 , \\
  \nabla \times \bm E &= -\frac{\partial \bm B}{\partial t} , \\
  \nabla \times \bm B &= \mu_0 \bm J + \mu_0 \varepsilon_0 \frac{\partial \bm E}{\partial t} .
\end{align*}
Integral form (Gauss and Ampère as examples):
\begin{align*}
  \oint_S \bm E \cdot d\bm S &= \frac{Q(W)}{\varepsilon_0} , \\
  \oint_C \bm B \cdot d\bm l &= \mu_0 I(S) + \mu_0 \varepsilon_0 \frac{d}{dt} \int_S \bm E \cdot d\bm S .
\end{align*}
Space is smooth and three-dimensional at all scales.
\end{frame}


%------------------------------------------------
\begin{frame}{Where to add Fractional Calculus?}
\begin{itemize}
  \item So far, no fractional \emph{derivatives} appeared in Maxwell equations
  \item If we naively replaced $\partial \rightarrow \partial^{\alpha}$ we would break Lorentz invariance
  \item But we can use Integration of non-integer order
  \item This corresponds to \emph{fractional measures} on $\mathbb{R}^3$
\end{itemize}

\vspace{0.2cm}
This approach follows Tarasov’s
\[
\text{``fractional continuous model of fractal distributions.''}
\]
\end{frame}


%------------------------------------------------
\begin{frame}{Fractional Integrals and Fractal Measures}
For a function $f(\bm r)$, the Riesz fractional integral in $\mathbb{R}^3$ is
\begin{equation*}
  (I^D f)(\bm r)
  =
  C(D)
  \int_{\mathbb{R}^3}
  \frac{f(\bm r')}{|\bm r - \bm r'|^{3-D}}
  \, d^3 r' ,
  \qquad 0 < D < 3
\end{equation*}

\begin{itemize}
  \item This integral has the same scaling as a $D$-dimensional volume
  \item It defines an \textbf{effective integration over a fractal}
  \item This is important since many fractals are singular sets! (Lebeasgue measure is 0)
\end{itemize}

\vspace{0.2cm}
Tarasov shows that
\[
  (I^D f)(\bm r)
  \;\;\Longleftrightarrow\;\;
  \int_W f(\bm r)\, dV_D
\]
\end{frame}



%------------------------------------------------
\begin{frame}{Density of States as Fractional Kernel}
Let's introduce the fractional volume element
\begin{equation*}
  dV_D = c_3(D,\bm r)\, d^3x
  \qquad
  c_3(D,r) \propto r^{D-3}
\end{equation*}

It is exactly the kernel of a Riesz-type fractional integral


\begin{itemize}
  \item Fractional integration replaces
  \[
    \int d^3x \;\;\to\;\; \int r^{D-3} d^3x
  \]
  \item The order of integration $D$ equals the fractal dimension
  \item Replace the fractal by an \textbf{effective continuous distribution} in $\mathbb{R}^3$
\end{itemize}

For a ball of radius $R$ one has
\begin{equation*}
  \int_{|\bm r| \le R} dV_D \propto R^D
\end{equation*}

\end{frame}


%------------------------------------------------
\begin{frame}{Density of States $c_n(D,\bm r)$}
\begin{itemize}
  \item $c_n(D,\bm r)$: \textbf{density of states} in $n$-dimensional Euclidean space
  \item Tells us how densely “available positions” for charges are packed
  \item For $n=3$, a typical choice (Riemann--Liouville type) is
    \begin{equation*}
      c_3(D,r) = C(D)\, r^{D-3}
    \end{equation*}
  where $r = |\bm r|$ and $C(D)$ is a dimension-dependent normalization factor
  \item For $D = 3$, $c_3(3,r) = 1$ and we recover the usual volume element
  \item Similar definitions exist for $c_2(d,\bm r)$ on surfaces and $c_1(\gamma,\bm r)$ on curves
\end{itemize}
\end{frame}

%------------------------------------------------
\begin{frame}{Generalized Gauss and Stokes Theorems}
\begin{itemize}
  \item For a vector field $\bm F(\bm r)$, Gauss theorem in this model is schematically
    \begin{equation*}
      \oint_{\partial W} \bm F \cdot d\bm S_d
      = \int_W \nabla \cdot \bigl( c_2(d,\bm r)\, \bm F(\bm r) \bigr)\, c_3(D,\bm r)\, d^3x
    \end{equation*}
  \item Stokes theorem similarly becomes
    \begin{equation*}
      \oint_{\partial S} \bm F \cdot d\bm l_\gamma
      = \int_S \bigl[ \nabla \times \bigl( c_1(\gamma,\bm r)\, \bm F(\bm r) \bigr) \bigr] \cdot d\bm S_d
    \end{equation*}
  \item The \emph{form} of the theorems is preserved, but the divergence and curl are effectively weighted by $c_n$.
\end{itemize}
\end{frame}

%------------------------------------------------
\begin{frame}{Charge and Current in Fractal Media}
\begin{itemize}
  \item Total fractal charge in region $W$:
    \begin{equation*}
      Q_D(W,t) = \int_W \rho(\bm r,t)\, c_3(D,\bm r)\, d^3x
    \end{equation*}
  \item Current density remains
    \begin{equation*}
      \bm J(\bm r,t) = \rho(\bm r,t)\, \bm u(\bm r,t)
    \end{equation*}
  \item but the total current through a fractal surface $S$ of dimension $d$ is
    \begin{equation*}
      I_d(S,t) = \int_S \bm J(\bm r,t) \cdot d\bm S_d, \qquad
      d\bm S_d = c_2(d,\bm r)\, d\bm S
    \end{equation*}
\end{itemize}
\end{frame}

%------------------------------------------------
\begin{frame}{Charge Conservation}
Global charge conservation for a region $W$:
\begin{equation*}
  \frac{d}{dt} Q_D(W,t) = - I_d(\partial W,t) .
\end{equation*}
In terms of densities:
\begin{equation*}
  \frac{d}{dt} \int_W \rho(\bm r,t)\, c_3(D,\bm r)\, d^3x
  = - \int_{\partial W} \bm J(\bm r,t) \cdot d\bm S_d
\end{equation*}
Using generalized Gauss theorem and assuming $W$ is fixed in time, this leads to a continuity equation:
\begin{equation*}
  c_3(D,\bm r)\, \frac{\partial \rho}{\partial t}
  + \nabla \cdot \bigl( c_2(d,\bm r)\, \bm J(\bm r,t) \bigr) = 0
\end{equation*}
For $D = 3$ and $d = 2$, $c_3 = c_2 = 1$ and we recover
\begin{equation*}
  \frac{\partial \rho}{\partial t} + \nabla \cdot \bm J = 0
\end{equation*}
\end{frame}

%------------------------------------------------
\begin{frame}{Coulomb’s Law for Fractal Sources}
For a charge distribution $\rho(\bm r')$ with fractal dimension $D$, the electric field at $\bm r$ is
\begin{equation*}
  \bm E(\bm r) = \frac{1}{4\pi \varepsilon_0}
  \int_W \frac{\bm r - \bm r'}{\lvert \bm r - \bm r' \rvert^3}\,
  \rho(\bm r')\, c_3(D,\bm r')\, d^3x'
\end{equation*}
\begin{itemize}
  \item The kernel is the same as in standard 3D electrostatics
  \item The \emph{measure} $c_3(D,\bm r')\, d^3x'$ carries information about fractal geometry
  \item For homogeneous fractal distributions, one obtains field scalings like $E(R) \propto R^{D-3}$
\end{itemize}
\end{frame}

%------------------------------------------------
\begin{frame}{Biot--Savart Law for Fractal Currents}
For a steady current distribution $\bm J(\bm r')$ with fractal dimension $D$:
\begin{equation*}
  \bm B(\bm r) = \frac{\mu_0}{4\pi}
  \int_W \frac{\bm J(\bm r') \times (\bm r - \bm r')}{\lvert \bm r - \bm r' \rvert^3}\,
  c_3(D,\bm r')\, d^3x'
\end{equation*}
\begin{itemize}
  \item Again, the kernel is unchanged
  \item Fractal integration modifies how strongly regions contribute as a function of distance
  \item In symmetric cases, one can get magnetic fields that decay with non-standard exponents or even stay constant
\end{itemize}
\end{frame}

%------------------------------------------------
\begin{frame}{Gauss and Amp\`ere Laws with Fractal Measures}
Electric flux through a fractal surface $S = \partial W$:
\begin{equation*}
  \Phi_E(S,t) = \oint_S \bm E(\bm r,t) \cdot d\bm S_d
  = \frac{1}{\varepsilon_0} \int_W \rho(\bm r,t)\, dV_D
\end{equation*}
Amp\`ere’s law (without displacement current, for steady currents):
\begin{equation*}
  \oint_C \bm B(\bm r) \cdot d\bm l_\gamma
  = \mu_0 \int_S \bm J(\bm r)\cdot d\bm S_d
\end{equation*}
\begin{itemize}
  \item For homogeneous fractal distributions, these lead to simple power-law fields
  \item Example: for suitable $D$ and $d$, $E(R)$ or $B(R)$ can be constant with $R$
\end{itemize}
\end{frame}

%------------------------------------------------
\begin{frame}{Maxwell Equations for Fractal Distributions (Integral Form)}
Assuming charges and fields are defined on the fractal set:
\begin{align*}
  \oint_S \bm E \cdot d\bm S_d
    &= \frac{1}{\varepsilon_0} \int_W \rho(\bm r,t)\, dV_D \\
  \oint_S \bm B \cdot d\bm S_d
    &= 0 \\
  \oint_C \bm E \cdot d\bm l_\gamma
    &= -\frac{d}{dt} \oint_S \bm B \cdot d\bm S_d \\
  \oint_C \bm B \cdot d\bm l_\gamma
    &= \mu_0 \int_S \bm J(\bm r,t)\cdot d\bm S_d
       + \mu_0 \varepsilon_0 \frac{d}{dt} \oint_S \bm E \cdot d\bm S_d
\end{align*}
\begin{itemize}
  \item Same structure as standard Maxwell equations, but with $dV_D$, $d\bm S_d$, $d\bm l_\gamma$.
\end{itemize}
\end{frame}

%------------------------------------------------
\begin{frame}{Maxwell Equations for Fractal Distributions (Differential Form)}
Using generalized Gauss and Stokes theorems, the integral equations lead to:
\begin{align*}
  \nabla \cdot \bigl( c_2(d,\bm r)\, \bm E(\bm r,t) \bigr)
    &= \frac{1}{\varepsilon_0} c_3(D,\bm r)\, \rho(\bm r,t) \\
  \nabla \cdot \bigl( c_2(d,\bm r)\, \bm B(\bm r,t) \bigr)
    &= 0 \\
  \nabla \times \bigl( c_1(\gamma,\bm r)\, \bm E(\bm r,t) \bigr)
    &= - c_2(d,\bm r)\, \frac{\partial \bm B}{\partial t} \\
  \nabla \times \bigl( c_1(\gamma,\bm r)\, \bm B(\bm r,t) \bigr)
    &= \mu_0 c_2(d,\bm r)\, \bm J(\bm r,t)
     + \mu_0 \varepsilon_0 c_2(d,\bm r)\, \frac{\partial \bm E}{\partial t}
\end{align*}
\begin{itemize}
  \item These are ordinary differential equations in space and time, but with position-dependent weights from geometry.
\end{itemize}
\end{frame}

%------------------------------------------------
\begin{frame}{Apparent Magnetic Monopoles from Fractality}
From
\begin{equation*}
  \nabla \cdot \bigl( c_2(d,\bm r)\, \bm B \bigr) = 0
\end{equation*}
we can write
\begin{equation*}
  c_2(d,\bm r)\, \nabla \cdot \bm B + \bm B \cdot \nabla c_2(d,\bm r) = 0
\end{equation*}
If $c_2$ varies in space (true fractal case), then
\begin{equation*}
  \nabla \cdot \bm B = - \frac{\bm B \cdot \nabla c_2}{c_2} \neq 0
\end{equation*}
in general.
\begin{itemize}
  \item The field behaves \emph{as if} there were effective magnetic charges
  \item These are not fundamental monopoles, but artifacts of the fractal geometry
\end{itemize}
\end{frame}

%------------------------------------------------
\begin{frame}{Effective Medium Interpretation}
Define effective fields
\begin{equation*}
  \bm E_{\text{eff}} = c_1(\gamma,\bm r)\, \bm E, \qquad
  \bm B_{\text{eff}} = c_2(d,\bm r)\, \bm B
\end{equation*}
Then we can rewrite Maxwell equations in a more standard form using
\begin{equation*}
  \bm D = \varepsilon_{\text{eff}}(\bm r)\, \bm E_{\text{eff}}, \qquad
  \bm H = \frac{1}{\mu_{\text{eff}}(\bm r)}\, \bm B_{\text{eff}}
\end{equation*}
\begin{itemize}
  \item Fractal distribution behaves like a medium with spatially varying
  \begin{equation*}
    \varepsilon_{\text{eff}}(\bm r), \qquad \mu_{\text{eff}}(\bm r),
  \end{equation*}
  determined by $c_n(D,\bm r)$ and the fractal dimensions.
  \item Geometry (fractal) and material response become deeply intertwined.
\end{itemize}
\end{frame}

%------------------------------------------------
\begin{frame}{Electric Multipole Expansion on Fractals}
Scalar potential for a static fractal charge distribution:
\begin{equation*}
  \phi(\bm R) =
  \frac{1}{4\pi \varepsilon_0}
  \int_W \frac{\rho(\bm r)\, dV_D}{\lvert \bm R - \bm r \rvert}
\end{equation*}
For $R = |\bm R|$ large compared to $|\bm r|$, expand
\begin{equation*}
  \frac{1}{\lvert \bm R - \bm r \rvert}
  = \sum_{n=0}^\infty \frac{r^n}{R^{n+1}} P_n(\cos\theta),
\end{equation*}
where $\theta$ is the angle between $\bm R$ and $\bm r$.
\begin{itemize}
  \item Leads to usual monopole, dipole, quadrupole, \dots\ terms
  \item Coefficients are now integrals with respect to $dV_D$
\end{itemize}
\end{frame}

%------------------------------------------------
\begin{frame}{Monopole and Dipole Moments}
Monopole term (total charge):
\begin{equation*}
  \phi_0(\bm R) = \frac{1}{4\pi \varepsilon_0} \frac{Q_D(W)}{R}, \qquad
  Q_D(W) = \int_W \rho(\bm r)\, dV_D
\end{equation*}
Dipole moment for fractal distribution:
\begin{equation*}
  \bm p^{(D)} = \int_W \bm r\, \rho(\bm r)\, dV_D
\end{equation*}
Dipole contribution:
\begin{equation*}
  \phi_1(\bm R) = \frac{1}{4\pi \varepsilon_0}
  \frac{\bm p^{(D)} \cdot \bm R}{R^3}
\end{equation*}
\begin{itemize}
  \item Structure identical to standard EM
  \item Numerical values change because $dV_D$ redistributes weight in space
  \item In simple homogeneous examples, deviations from the $D=3$ dipole are typically of order $10\text{--}20\%$.
\end{itemize}
\end{frame}

%------------------------------------------------
\begin{frame}{Comparison with Standard Electrodynamics}
\begin{center}
\begin{tabular}{l|l}
\textbf{Standard EM} & \textbf{Fractal EM} \\\hline
Space dimension $3$ & Effective dimension $D$ (non-integer) \\
Ordinary measure $d^3x$ & Fractional measure $dV_D = c_3(D,\bm r)\, d^3x$ \\
Uniform media & Geometrically inhomogeneous media \\
$\nabla \cdot \bm B = 0$ & $\nabla \cdot (c_2 \bm B) = 0$ (effective monopoles) \\
Field scaling $E \sim 1/R^2$ & $E \sim R^{D-3}$ (in simple symmetric cases)
\end{tabular}
\end{center}
\end{frame}

%------------------------------------------------
\begin{frame}{Magnetohydrodynamics on Fractal Media}
\begin{itemize}
  \item Consider a plasma or conducting fluid whose \emph{charged particles}
        are distributed on a fractal set of dimension $D$
  \item The electromagnetic field interacts with a \textbf{collective medium},
        not isolated charges
\end{itemize}

\textbf{Mass (or charge) conservation:}
\begin{equation*}
  c_3(D,\bm r)\, \frac{\partial \rho}{\partial t}
  + \nabla \cdot \bigl( c_2(d,\bm r)\, \rho\, \bm u \bigr) = 0
\end{equation*}

\vspace{0.2cm}
\textbf{Momentum balance (Euler equation):}
\begin{equation*}
  c_3(D,\bm r)\, \rho
  \left(
    \frac{\partial \bm u}{\partial t}
    + (\bm u \cdot \nabla)\bm u
  \right)
  = - c_3(D,\bm r)\, \nabla p
    + \rho\, \bm E
    + \bm J \times \bm B
\end{equation*}

\begin{itemize}
  \item Same structure as standard fluid dynamics
  \item Inertia and pressure are weighted by the fractal measure
  \item Lorentz force keeps its usual form
\end{itemize}
\end{frame}


%------------------------------------------------
\begin{frame}{Induction Equation for Fractal Conducting Media}
Using generalized Maxwell equations and Ohm's law
\[
  \bm J = \sigma \left( \bm E + \bm u \times \bm B \right),
\]
where $\sigma$ is conductivity, one obtains the \textbf{fractal induction equation}:
\begin{equation*}
  \frac{\partial}{\partial t}
  \bigl( c_2(d,\bm r)\, \bm B \bigr)
  = \nabla \times
  \Bigl[
    c_1(\gamma,\bm r)\,
    \bigl( \bm u \times \bm B \bigr)
  \Bigr]
  +
  \eta\, \nabla^2
  \bigl( c_2(d,\bm r)\, \bm B \bigr)
\end{equation*}
where $\eta = 1/(\mu_0 \sigma)$ is the magnetic diffusivity.

\begin{itemize}
  \item Geometry modifies magnetic flux conservation
  \item Even ideal MHD ($\eta \to 0$) is affected by fractality
  \item Flux freezing becomes \emph{scale-dependent}
\end{itemize}
\end{frame}


%------------------------------------------------
\begin{frame}{Physical Interpretation and Limits}
\begin{itemize}
  \item Fractal geometry acts as an \textbf{effective medium}
  \item Transport, diffusion, and field amplification depend on $D$, $d$, $\gamma$
  \item Standard MHD is recovered for
  \[
    D = 3, \quad d = 2, \quad \gamma = 1
    \;\Rightarrow\;
    c_n = 1
  \]
  \item Deviations become important in:
  \begin{itemize}
    \item Turbulent plasmas
    \item Percolating current networks
    \item Space and astrophysical plasmas
  \end{itemize}
\end{itemize}
\end{frame}


%------------------------------------------------
\begin{frame}
  \centering
  {\Huge Thank you!}
\end{frame}



\end{document}